\begin{abstract}
本文研究移动机器人在部分可观测环境中对多个无线电干扰源进行主动搜索、定位与清除的问题。围绕测向误差、未知接收半径、多频道切换以及定向发射等因素，建立由几何定位、信念更新和路径决策组成的统一模型。

\textbf{针对问题一，}将每次测向观测表示为带角度误差的扇形可行域，通过半平面裁剪求取多次观测的交集，并利用最远点对和最小包围圆刻画定位精度与清除可行性。\placeholder{补充最终精度、误差或判定阈值。}

\textbf{针对问题二，}建立以单位虚拟时间信息增益为指标的下一最佳观测点模型，兼顾交会角、移动距离和失去信号风险。\placeholder{补充第二检测点坐标及相较基线的提升。}

\textbf{针对问题三，}把多源搜索抽象为部分可观测决策过程，以粒子集合维护各频道的位置后验，结合候选动作生成与有限步信念空间规划，联合优化测量、移动、清除和频道切换。\placeholder{补充正式模拟器中的清除率和平均耗时。}

\textbf{针对问题四，}进一步引入未知发射方向与无信号观测的多义性，通过定向源状态扩维、域随机化和覆盖兜底策略提高模型稳健性。\placeholder{补充问题四的核心数值结论。}

最后，从定位误差、策略稳定性和计算效率三个方面检验模型。结果表明，\placeholder{用一句话填写最有竞争力、且已经由程序复现的结论。}

\keywords{主动定位\quad 信念空间规划\quad 部分可观测决策过程\quad 信息增益\quad 路径优化}
\end{abstract}
